% Generated by experiments/make_tables.py from experiments/results/e1_re1/summary_by_system_fault.csv.
% Do not edit by hand.
\begin{table}[t]
\centering
\caption{Table 2b. E1, secondary: loose ExpAC@1 per fault type}
\label{tab:e1_loose}
\small
\begin{tabular}{lrrrrr}
\toprule
System / fault & Acceptable metrics & Isolation Forest + TreeSHAP & PCA reconstruction error & ECDF tail score (ECOD-style) & Robust z-score (PARE without residuals) \\
\midrule
\multicolumn{6}{l}{\textit{Online Boutique}} \\
CPU & 1 & 0.120 & 0.480 & 0.314 & 0.080 \\
MEM & 1 & 0.360 & 0.120 & 0.170 & 0.320 \\
DISK & 4 & 0.400 & 0.680 & 0.081 & 0.360 \\
DELAY & 3 & 0.320 & 1.000 & 0.122 & 0.480 \\
LOSS & 3 & 0.400 & 0.320 & 0.117 & 0.200 \\
\multicolumn{6}{l}{\textit{Sock Shop}} \\
CPU & 1 & 0.000 & 0.040 & 0.056 & 0.040 \\
MEM & 1 & 0.040 & 0.280 & 0.043 & 0.240 \\
DISK & 3 & 0.200 & 0.760 & 0.074 & 0.320 \\
DELAY & 2 & 0.400 & 0.840 & 0.055 & 0.480 \\
LOSS & 2 & 0.440 & 0.600 & 0.068 & 0.320 \\
\multicolumn{6}{l}{\textit{Train Ticket}} \\
CPU & 1 & 0.040 & 0.320 & 0.027 & 0.000 \\
MEM & 1 & 0.040 & 0.880 & 0.021 & 0.280 \\
DISK & 3 & 0.000 & 0.080 & 0.029 & 0.000 \\
DELAY & 2 & 0.120 & 0.560 & 0.028 & 0.080 \\
LOSS & 2 & 0.120 & 0.280 & 0.028 & 0.040 \\
\bottomrule
\end{tabular}
\par\smallskip\footnotesize
A hit when the top indicator is any metric in the fault type's acceptable set; the set's size is in the second column, and a larger set is an easier question.\par
Reported per fault type only, never pooled: averaging a 1-metric question with a 4-metric one would mean nothing (docs/08 section 4).\par
For CPU and MEM the set holds one metric, so these rows equal the strict score.\par
\end{table}
